% Source: Gerzensee "Calculus" slides 55-64 (the least-squares problem; Calculus
% Approach I with the two normal equations and the explicit 2x2 inverse; Calculus
% Approach II with the matrix form X^T y = X^T X beta; the projection aside, the
% orthogonality picture with 1 and x spanning the fitted subspace, and the
% projection matrix X(X^T X)^{-1} X^T; Example 8.1, the n=500 simulation).
% Supplementary: Fukuda (2026), Ch. 9 -- Thm 9.1 and Prop. 9.1 (projection onto a
% closed convex set), Def. 9.2-9.3 and Prop. 9.4-9.6 (orthogonality),
% Prop. 9.8-9.9 and Thm 9.2 (projection onto a closed subspace), Thm 9.3 (Riesz),
% Thm 9.4 (projection in R^n), Sec. 9.4 (OLS three ways), Sec. 9.5 (Gram-Schmidt
% and QR), Sec. 9.6.2 (best linear predictor).
% Code: Supplementary/codes/regression.m, regression_Octave.m,
% regression_Python.py (the Example 8.1 simulation).
% No external reference is cited in this chapter: every statement is drawn
% from the Gerzensee slides and the supplementary lecture notes listed above.

\chapter{Projection, Orthogonality, and Least Squares}
\label{ch:projection}

\section{Motivation: Least Squares Is Not About Calculus}\label{sec:proj-motivation}

\autoref{prop:ols-calculus} solved the least-squares problem by setting a
gradient to zero. The solution was correct and the method was, in a precise
sense, beside the point. Writing $\mathbf1=(1,\dots,1)^{\transp}$,
$x=(x_1,\dots,x_n)^{\transp}$ and $y=(y_1,\dots,y_n)^{\transp}$ in $\R^n$, the
objective is
\begin{equation}\label{eq:ols-norm}
	S(\beta_1,\beta_2)=\sum_{i=1}^n\bigl(y_i-\beta_1-\beta_2x_i\bigr)^2
	=\norm{y-\beta_1\mathbf1-\beta_2x}^2 ,
\end{equation}
so the problem is to find the point of the plane $\Span\{\mathbf1,x\}$ closest to
$y$ \autocite[slide 61]{fukuda2026calculus}. That is a question about distance in
an inner product space, and its answer --- drop a perpendicular --- does not
involve differentiation, holds in infinite dimensions where the calculus argument
would not, and explains why the first-order conditions look the way they do:
the normal equations \emph{are} the statement that the residual is orthogonal to
the regressors.

This chapter develops that geometry. Its three payoffs are the orthogonal
decomposition theorem, the projection formula
$P=X(X^{\transp}X)^{-1}X^{\transp}$, and the identification of conditional
expectation as a projection, which \autoref{ch:probability} uses to characterise
the best predictor.

\section{Orthogonality}\label{sec:proj-orthogonality}

Throughout, $H$ is a real inner product space with inner product $\ip{\cdot}{\cdot}$
and induced norm $\norm{\cdot}$ (\autoref{prop:ip-induces-norm}), and it is a
Hilbert space when complete (\autoref{def:banach-hilbert}). The leading cases are
$H=\R^n$ and $H=L^2$.

\begin{definition}[Orthogonal complement]\label{def:orthogonal-complement}
	For $\emptyset\neq A\subseteq H$, the \dfn{orthogonal complement} is
	\[
		A^{\perp}\eqdef\{x\in H : \ip{x}{a}=0\ \text{for all }a\in A\} .
	\]
	A set $A$ is \dfnas{orthogonal}{orthogonal set} if $\ip{x}{y}=0$ for all
	distinct $x,y\in A$, and \dfnas{orthonormal}{orthonormal set} if in addition
	$\norm{x}=1$ for every $x\in A$.
	\nidx{Aperp}{$A^{\perp}$, orthogonal complement}
\end{definition}

\begin{proposition}[Elementary properties]\label{prop:orthogonal-complement}
	Let $A,B\subseteq H$ be non-empty.
	\begin{enumerate}[(i)]
		\item $A^{\perp}$ is a closed vector subspace of $H$, whether or not $A$ is.
		\item $A\subseteq B$ implies $B^{\perp}\subseteq A^{\perp}$, and
		      $A\subseteq A^{\perp\perp}$, and $A^{\perp}=A^{\perp\perp\perp}$.
		\item $x\in\bigl(\Span\{x_1,\dots,x_m\}\bigr)^{\perp}$ if and only if
		      $\ip{x}{x_i}=0$ for every $i$.
		\item An orthogonal set not containing $0$ is linearly independent.
	\end{enumerate}
\end{proposition}

\begin{proof}
	(i) Linearity of the inner product in its first argument makes $A^{\perp}$ a
	subspace. For closedness, let $x_k\to x$ with $x_k\in A^{\perp}$; for each
	$a\in A$, $\ip{x}{a}=\lim_k\ip{x_k}{a}=0$ by continuity of the inner product,
	which follows from Cauchy--Schwarz (\autoref{lem:cauchy-schwarz}).

	(ii) The first claim is immediate. For $A\subseteq A^{\perp\perp}$: if $a\in A$
	then $\ip{a}{z}=0$ for every $z\in A^{\perp}$ by definition of $A^{\perp}$, so
	$a\in A^{\perp\perp}$. Applying the first claim to
	$A\subseteq A^{\perp\perp}$ gives $A^{\perp\perp\perp}\subseteq A^{\perp}$, and
	applying $A\subseteq A^{\perp\perp}$ with $A^{\perp}$ in place of $A$ gives the
	reverse inclusion.

	(iii) If $\ip{x}{x_i}=0$ for every $i$, then $\ip{x}{\sum_i\lambda_ix_i}=0$ by
	linearity. The converse is the special case of singleton combinations.

	(iv) Let $x_1,\dots,x_m$ be orthogonal and non-zero with
	$\sum_i\lambda_ix_i=0$. Taking the inner product with $x_j$ kills every term but
	the $j$th, giving $\lambda_j\norm{x_j}^2=0$, hence $\lambda_j=0$.
\end{proof}

These are Propositions~9.4 and~9.6 of \textcite[pp.~299--301]{fukuda2026notes}.
Part (i) contains a fact used repeatedly: orthogonal complements are
automatically closed, which is why the decomposition theorem below can be stated
for an arbitrary closed subspace $M$ without any hypothesis on $M^{\perp}$.

\begin{proposition}[Pythagoras]\label{prop:pythagoras}
	If $x\perp y$ then $\norm{x+y}^2=\norm{x}^2+\norm{y}^2$. More generally, for a
	finite orthogonal set $\{x_1,\dots,x_m\}$,
	$\norm{\sum_ix_i}^2=\sum_i\norm{x_i}^2$.
\end{proposition}

\begin{proof}
	Expand $\norm{x+y}^2=\norm{x}^2+2\ip{x}{y}+\norm{y}^2$ and use $\ip{x}{y}=0$;
	iterate for the general case.
\end{proof}

\section{The Projection Theorem}\label{sec:proj-theorem}

\begin{theorem}[Projection onto a closed convex set]\label{thm:projection-convex}
	Let $H$ be a Hilbert space, $\emptyset\neq M\subseteq H$ closed and convex, and
	$x\in H$. Then there is a unique $P_Mx\in M$ with
	\[
		\norm{x-P_Mx}=\inf_{y\in M}\norm{x-y},
	\]
	and $y_0=P_Mx$ is characterised by
	\begin{equation}\label{eq:proj-variational}
		\ip{x-y_0}{z-y_0}\leq0\qquad\text{for all }z\in M .
	\end{equation}
\end{theorem}

\begin{proof}
	\emph{Existence.} Let $d\eqdef\inf_{y\in M}\norm{x-y}$ and take
	$(y_k)\subseteq M$ with $\norm{x-y_k}\to d$. The parallelogram law
	(\autoref{rem:parallelogram}) applied to $x-y_j$ and $x-y_k$ gives
	\[
		\norm{y_j-y_k}^2
		=2\norm{x-y_j}^2+2\norm{x-y_k}^2-4\left\lVert x-\frac{y_j+y_k}{2}\right\rVert^2
		\leq2\norm{x-y_j}^2+2\norm{x-y_k}^2-4d^2,
	\]
	where the inequality uses $\tfrac12(y_j+y_k)\in M$ by convexity. As
	$j,k\to\infty$ the right side tends to $2d^2+2d^2-4d^2=0$, so $(y_k)$ is
	Cauchy. Completeness of $H$ and closedness of $M$
	(\autoref{prop:closed-subset-complete}) give a limit $y_0\in M$, and
	$\norm{x-y_0}=\lim_k\norm{x-y_k}=d$ by continuity of the norm.

	\emph{Uniqueness and characterisation.} Identical to the corresponding parts of
	the proof of \autoref{lem:nearest-point}, which used only the parallelogram law
	and convexity.
\end{proof}

This is Theorem~9.1 with Proposition~9.1 of
\textcite[pp.~295--298]{fukuda2026notes}. The comparison with
\autoref{lem:nearest-point} is instructive. There, existence came from
compactness: a closed bounded subset of $\R^n$ is compact and a continuous
function attains its minimum. Here compactness is unavailable --- the closed unit
ball of an infinite-dimensional space is not compact
(\autoref{rem:heine-borel-fails}) --- and the parallelogram law does the work
instead, converting a minimising sequence into a Cauchy sequence, which
completeness then converges. Convexity is used exactly once, to place the
midpoint $\tfrac12(y_j+y_k)$ in $M$, and the theorem is false without it: in
$\R^2$ the point $0$ has no unique nearest point in the closed non-convex set
$\{y:\norm{y}=1\}$.

\begin{proposition}[Characterisation for subspaces]\label{prop:proj-subspace-char}
	Let $M\subseteq H$ be a closed subspace and $x\in H$. Then $y_0\in M$ minimises
	$\norm{x-y}$ over $M$ if and only if
	\begin{equation}\label{eq:orthogonality-condition}
		\ip{x-y_0}{z}=0\qquad\text{for all }z\in M,
	\end{equation}
	that is, if and only if $x-y_0\in M^{\perp}$.
\end{proposition}

\begin{proof}
	($\Rightarrow$) Let $u\in M$ and $t\in\R$. Then $y_0+tu\in M$, so
	\[
		\norm{x-y_0}^2\leq\norm{x-y_0-tu}^2
		=\norm{x-y_0}^2-2t\ip{x-y_0}{u}+t^2\norm{u}^2 .
	\]
	Hence $t^2\norm{u}^2-2t\ip{x-y_0}{u}\geq0$ for \emph{every} real $t$,
	including $t$ of both signs. A quadratic in $t$ with a non-negative value
	everywhere and a zero at $t=0$ must have vanishing linear coefficient, so
	$\ip{x-y_0}{u}=0$.

	($\Leftarrow$) If $x-y_0\perp M$ then for any $y\in M$, $y_0-y\in M$ and
	Pythagoras (\autoref{prop:pythagoras}) gives
	\[
		\norm{x-y}^2=\norm{(x-y_0)+(y_0-y)}^2=\norm{x-y_0}^2+\norm{y_0-y}^2
		\geq\norm{x-y_0}^2 .
	\]
\end{proof}

This is Proposition~9.8 of \textcite[p.~301]{fukuda2026notes}. Comparing
\eqref{eq:orthogonality-condition} with \eqref{eq:proj-variational}: on a
subspace the variational inequality holds for $z$ and for $-z$ simultaneously,
so it collapses to an equality. The last display also proves more than was asked:
it quantifies the loss from using any other $y\in M$, and that identity is the
decomposition of total variation used in \autoref{eq:anova} below.

\begin{theorem}[Orthogonal decomposition]\label{thm:orthogonal-decomposition}
	Let $M$ be a closed subspace of a Hilbert space $H$. Then every $x\in H$ has a
	unique decomposition
	\[
		x=y+z,\qquad y=P_Mx\in M,\quad z\in M^{\perp},
	\]
	so that $H=M\oplus M^{\perp}$, and $M^{\perp\perp}=M$.
\end{theorem}

\begin{proof}
	\emph{Existence.} If $x\in M$ take $(y,z)=(x,0)$. Otherwise let $y=P_Mx$ from
	\autoref{thm:projection-convex}, which applies since a closed subspace is closed
	and convex, and let $z\eqdef x-y$; then $z\in M^{\perp}$ by
	\autoref{prop:proj-subspace-char}.

	\emph{Uniqueness.} If $y_1+z_1=y_2+z_2$ with $y_i\in M$, $z_i\in M^{\perp}$,
	then $y_1-y_2=z_2-z_1\in M\cap M^{\perp}$. Any $w$ in that intersection
	satisfies $\ip{w}{w}=0$, hence $w=0$.

	\emph{$M^{\perp\perp}=M$.} The inclusion $M\subseteq M^{\perp\perp}$ is
	\autoref{prop:orthogonal-complement}(ii). Conversely, let
	$x\in M^{\perp\perp}$ and decompose $x=y+z$ as above. Then
	$z=x-y\in M^{\perp}$, and also $z\in M^{\perp\perp}$ since both $x$ and
	$y\in M\subseteq M^{\perp\perp}$ lie there. So $\ip{z}{z}=0$ and $x=y\in M$.
\end{proof}

This is Theorem~9.2 of \textcite[p.~302]{fukuda2026notes}. The hypothesis that
$M$ be \emph{closed} cannot be dropped, and in infinite dimensions it is not
automatic: by \autoref{cor:subspace-complete} every subspace of $\R^n$ is closed,
but the subspace of finitely supported sequences in $L^2$ is dense and proper, so
its orthogonal complement is $\{0\}$ and no decomposition of the stated form
exists.

\begin{proposition}[The projection operator]\label{prop:projection-operator}
	Let $M\subseteq H$ be a closed subspace. The map $P_M\colon H\to M$ is
	\begin{enumerate}[(i)]
		\item linear and surjective onto $M$;
		\item idempotent, $P_M^2=P_M$, with $\Ker P_M=M^{\perp}$ and
		      $\Image P_M=M$;
		\item self-adjoint, $\ip{P_Mx}{w}=\ip{x}{P_Mw}$ for all $x,w\in H$;
		\item norm-reducing, $\norm{P_Mx}\leq\norm{x}$, with equality if and only if
		      $x\in M$.
	\end{enumerate}
\end{proposition}

\begin{proof}
	(i) Let $x=y+z$ and $x'=y'+z'$ be the decompositions of
	\autoref{thm:orthogonal-decomposition} and $\alpha,\beta\in\R$. Then
	$\alpha x+\beta x'=(\alpha y+\beta y')+(\alpha z+\beta z')$ with the first
	bracket in $M$ and the second in $M^{\perp}$, both being subspaces; by
	uniqueness this is \emph{the} decomposition, so
	$P_M(\alpha x+\beta x')=\alpha P_Mx+\beta P_Mx'$. Surjectivity: $P_My=y$ for
	$y\in M$.

	(ii) Immediate from $P_My=y$ on $M$ and $P_Mz=0$ on $M^{\perp}$.

	(iii) Write $x=y+z$, $w=y'+z'$. Then
	$\ip{P_Mx}{w}=\ip{y}{y'+z'}=\ip{y}{y'}$ since $y\in M$ and $z'\in M^{\perp}$,
	and symmetrically $\ip{x}{P_Mw}=\ip{y+z}{y'}=\ip{y}{y'}$.

	(iv) $\norm{x}^2=\norm{y}^2+\norm{z}^2\geq\norm{y}^2$ by
	\autoref{prop:pythagoras}, with equality exactly when $z=0$.
\end{proof}

This is Proposition~9.9 of \textcite[p.~302]{fukuda2026notes}, with the
self-adjointness added because it is what identifies the projection matrix in
\autoref{prop:projection-matrix} and what makes the residual maker symmetric.

\begin{theorem}[Riesz representation]\label{thm:riesz-hilbert}
	Let $H$ be a Hilbert space and $f\colon H\to\R$ a continuous linear functional.
	Then there is a unique $a\in H$ with $f(x)=\ip{a}{x}$ for all $x\in H$, and
	$\norm{a}$ equals the operator norm of $f$.
\end{theorem}

\begin{proof}
	If $f\equiv0$ take $a=0$. Otherwise $M\eqdef\Ker f$ is a closed proper subspace
	--- closed because $f$ is continuous and $M=f^{-1}(\{0\})$
	(\autoref{thm:continuity-preimage}) --- so by
	\autoref{thm:orthogonal-decomposition} there is $u\in M^{\perp}$ with
	$\norm{u}=1$. For any $x\in H$ the vector $f(x)u-f(u)x$ lies in $M$, since $f$
	of it is $f(x)f(u)-f(u)f(x)=0$; taking the inner product with $u\in M^{\perp}$
	gives $f(x)-f(u)\ip{u}{x}=0$, that is $f(x)=\ip{f(u)u}{x}$. So $a=f(u)u$ works.
	Uniqueness: if $\ip{a}{x}=\ip{a'}{x}$ for all $x$, take $x=a-a'$.
\end{proof}

This is Theorem~9.3 of \textcite[p.~304]{fukuda2026notes} and the promised
generalisation of \autoref{thm:riesz-rn}; in $\R^n$ continuity was automatic by
\autoref{prop:linear-continuous}(ii), and in infinite dimensions it must be
assumed, since discontinuous linear functionals exist there.

\section{Projection in \texorpdfstring{$\R^n$}{Rn} and the Projection Matrix}
\label{sec:proj-rn}

\begin{proposition}[Projection matrix]\label{prop:projection-matrix}
	Let $X\in M_{n\times k}(\R)$ have linearly independent columns, and let
	$M\eqdef\Image X=\Span\{X_{\cdot1},\dots,X_{\cdot k}\}\subseteq\R^n$. Then
	$X^{\transp}X$ is invertible, and for every $y\in\R^n$,
	\begin{equation}\label{eq:projection-matrix}
		P_My=X\bigl(X^{\transp}X\bigr)^{-1}X^{\transp}y,
		\qquad
		P\eqdef X\bigl(X^{\transp}X\bigr)^{-1}X^{\transp} .
	\end{equation}
	The matrix $P$ is symmetric and idempotent with $\rank P=k$, and
	$Q\eqdef I_n-P$ is the projection onto $M^{\perp}$, symmetric and idempotent
	with $\rank Q=n-k$.
\end{proposition}

\begin{proof}
	\emph{Invertibility.} Suppose $X^{\transp}Xv=0$ for some $v\in\R^k$. Then
	$0=v^{\transp}X^{\transp}Xv=\norm{Xv}^2$, so $Xv=0$, and linear independence of
	the columns forces $v=0$. Hence $\Ker(X^{\transp}X)=\{0\}$ and
	\autoref{thm:rank-nullity} makes the square matrix $X^{\transp}X$ invertible.

	\emph{The formula.} Set $\hat y\eqdef X(X^{\transp}X)^{-1}X^{\transp}y$, which
	lies in $M$ because it is $X$ times a vector. By
	\autoref{prop:proj-subspace-char} it suffices to check
	$y-\hat y\perp M$, and by \autoref{prop:orthogonal-complement}(iii) it suffices
	to check orthogonality to each column of $X$, that is
	$X^{\transp}(y-\hat y)=0$:
	\[
		X^{\transp}y-X^{\transp}X\bigl(X^{\transp}X\bigr)^{-1}X^{\transp}y
		=X^{\transp}y-X^{\transp}y=0 .
	\]

	\emph{Properties.} Symmetry: $P^{\transp}=X\bigl((X^{\transp}X)^{-1}\bigr)^{\transp}
		X^{\transp}=P$, since $X^{\transp}X$ is symmetric and so is its inverse.
	Idempotence: $P^2=X(X^{\transp}X)^{-1}(X^{\transp}X)(X^{\transp}X)^{-1}X^{\transp}=P$.
	Rank: $\Image P=M$ has dimension $k$. The statements for $Q$ follow from
	\autoref{thm:orthogonal-decomposition}, $Q=I-P$ being the projection onto
	$M^{\perp}$, whose dimension is $n-k$ by
	\autoref{thm:rank-nullity}.
\end{proof}

This is Theorem~9.4 of \textcite[p.~305]{fukuda2026notes}, and it is the
formula displayed on \autocite[slide 62]{fukuda2026calculus}. Both hypotheses on
$X$ carry economic content: linear independence of the columns is the
no-perfect-multicollinearity condition of \autoref{rem:square-matters}, and
$k\leq n$ is the requirement of at least as many observations as parameters,
which is forced because $k$ independent vectors cannot live in $\R^n$ with $k>n$
by \autoref{thm:dimension-well-defined}.

\subsection{Gram--Schmidt and QR}

\begin{proposition}[Gram--Schmidt]\label{prop:gram-schmidt}
	Let $v_1,\dots,v_k\in H$ be linearly independent. Define recursively
	\[
		u_1\eqdef v_1,\qquad
		u_j\eqdef v_j-\sum_{i=1}^{j-1}\frac{\ip{v_j}{u_i}}{\norm{u_i}^2}\,u_i
		\quad(j=2,\dots,k),
		\qquad e_j\eqdef\frac{u_j}{\norm{u_j}} .
	\]
	Then $\{e_1,\dots,e_k\}$ is orthonormal and
	$\Span\{e_1,\dots,e_j\}=\Span\{v_1,\dots,v_j\}$ for every $j$.
\end{proposition}

\begin{proof}
	Induction on $j$. Each $u_j$ is $v_j$ minus its projection onto
	$\Span\{u_1,\dots,u_{j-1}\}$, expressed through the orthogonal basis by
	\autoref{prop:pythagoras}; by \autoref{prop:proj-subspace-char} the difference
	is orthogonal to that span. It is non-zero because $v_j\notin
		\Span\{v_1,\dots,v_{j-1}\}=\Span\{u_1,\dots,u_{j-1}\}$ by independence, so
	the normalisation is legitimate, and the span statement follows since the
	transformation is triangular with unit diagonal.
\end{proof}

\begin{corollary}[QR decomposition]\label{cor:qr}
	Every $X\in M_{n\times k}(\R)$ with independent columns factors as $X=QR$ with
	$Q\in M_{n\times k}(\R)$ having orthonormal columns and $R\in M_{k\times k}(\R)$
	upper triangular with positive diagonal. Then
	$P=QQ^{\transp}$ and the least-squares coefficient solves $R\hat\beta=Q^{\transp}y$.
\end{corollary}

\begin{proof}
	Apply \autoref{prop:gram-schmidt} to the columns of $X$ and collect the
	coefficients of $v_j$ in terms of $e_1,\dots,e_j$ into the $j$th column of $R$,
	which is upper triangular by construction. Then
	$X^{\transp}X=R^{\transp}Q^{\transp}QR=R^{\transp}R$ and
	\eqref{eq:projection-matrix} gives
	$P=QR(R^{\transp}R)^{-1}R^{\transp}Q^{\transp}=QQ^{\transp}$; substituting into
	$X^{\transp}X\hat\beta=X^{\transp}y$ gives
	$R^{\transp}R\hat\beta=R^{\transp}Q^{\transp}y$, and $R^{\transp}$ is invertible.
\end{proof}

This is Section~9.5 of \textcite[pp.~320--323]{fukuda2026notes}. The practical
point is numerical: forming $X^{\transp}X$ squares the condition number of $X$,
so the explicit inverse in \eqref{eq:projection-matrix} is a formula for
analysis and a poor algorithm, whereas solving the triangular system
$R\hat\beta=Q^{\transp}y$ is stable. Statistical software computes the QR
factorisation, not the inverse.

\section{Econometric Application: Least Squares}\label{sec:proj-ols}

\paragraph{Setting.} Observations $(x_i,y_i)_{i=1}^n$ are stacked as
$y\in\R^n$ and $X=[\mathbf1\ \ x]\in M_{n\times2}(\R)$. The least-squares
problem is \eqref{eq:ols-norm}.

\begin{proposition}[Least squares as projection]\label{prop:ols-projection}
	Suppose the columns of $X$ are linearly independent, equivalently that the
	$x_i$ are not all equal. Then \eqref{eq:ols-norm} has the unique solution
	\begin{equation}\label{eq:ols-matrix}
		\hat\beta=\bigl(X^{\transp}X\bigr)^{-1}X^{\transp}y,
	\end{equation}
	the fitted vector is $\hat y=X\hat\beta=Py$, the residual is
	$\hat e=y-\hat y=Qy$, and
	\begin{equation}\label{eq:normal-equations-geometric}
		X^{\transp}\hat e=0,
		\qquad\text{that is}\qquad
		\sum_{i=1}^n\hat e_i=0
		\quad\text{and}\quad\sum_{i=1}^nx_i\hat e_i=0 .
	\end{equation}
	Moreover
	\begin{equation}\label{eq:anova}
		\norm{y-\bar y\mathbf1}^2=\norm{\hat y-\bar y\mathbf1}^2+\norm{\hat e}^2 .
	\end{equation}
\end{proposition}

\begin{proof}
	Minimising \eqref{eq:ols-norm} is minimising $\norm{y-v}$ over
	$v\in M=\Image X$, so the solution is $\hat y=P_My$, unique by
	\autoref{thm:projection-convex}, and given by
	\eqref{eq:projection-matrix}. Since $X$ has independent columns, the
	representation $\hat y=X\hat\beta$ has a unique coefficient vector, namely
	\eqref{eq:ols-matrix}. Orthogonality
	\eqref{eq:normal-equations-geometric} is
	\autoref{prop:proj-subspace-char} combined with
	\autoref{prop:orthogonal-complement}(iii), the two scalar equations being
	orthogonality to $\mathbf1$ and to $x$.

	For \eqref{eq:anova}: the first equation of
	\eqref{eq:normal-equations-geometric} gives
	$\bar{\hat e}=0$, hence $\bar{\hat y}=\bar y$, so
	$\hat y-\bar y\mathbf1\in M$ and $\hat e\in M^{\perp}$ are orthogonal;
	$y-\bar y\mathbf1=(\hat y-\bar y\mathbf1)+\hat e$ and
	\autoref{prop:pythagoras} apply.
\end{proof}

\paragraph{What the geometry explains.} Four things that the calculus derivation
leaves opaque.

First, the normal equations \eqref{eq:normal-equations-geometric} are not
computational artefacts. They say the residual is orthogonal to every regressor,
which is the defining property of the projection, and they are the sample analogue
of the population moment conditions $\E[\hat e]=0$, $\E[x\hat e]=0$ that identify
the estimand. The first equation holds \emph{because} $\mathbf1$ is a regressor:
a regression without an intercept has residuals that need not sum to zero, and
then $\bar{\hat y}\neq\bar y$ and \eqref{eq:anova} fails.

Second, \eqref{eq:anova} is the analysis-of-variance decomposition, total sum of
squares equals explained plus residual, and it is Pythagoras' theorem. The
coefficient of determination $R^2=\norm{\hat y-\bar y\mathbf1}^2/
	\norm{y-\bar y\mathbf1}^2$ is therefore $\cos^2$ of the angle between the
demeaned data vector and the demeaned fit, which by
\autoref{eg:correlation-cosine} is the squared sample correlation between $y$ and
$\hat y$. That $R^2\in[0,1]$ is Cauchy--Schwarz, not a convention.

Third, adding a regressor cannot raise the residual sum of squares, because it
enlarges $M$ and \autoref{thm:projection-convex} minimises over a larger set. This
is a geometric fact about nested subspaces, which is why $R^2$ is
non-decreasing in the number of regressors and why model selection needs a
penalty rather than a fit criterion.

Fourth, the estimator is linear in $y$: $\hat\beta=(X^{\transp}X)^{-1}X^{\transp}y$
applies a fixed matrix to the data. Linearity is
\autoref{prop:projection-operator}(i), and it is the property that makes the
sampling distribution of $\hat\beta$ computable from that of $y$ --- normal if
$y$ is normal, by the linear-image property of the Gaussian family.

\begin{eg}[The simulation]\label{eg:ols-simulation}
	Example~8.1 of \textcite{fukuda2026notes}
	\autocite[slides 63--64]{fukuda2026calculus} draws $n=500$ independent standard
	normal regressors $x_i$ and errors $\varepsilon_i$, sets
	$y_i=1+3x_i+\varepsilon_i$, and computes $\hat\beta$ from
	\eqref{eq:ols-matrix}. The accompanying programs
	--- \path{regression.m} for \textsc{Matlab}, \path{regression_Octave.m}, and
	\path{regression_Python.py} --- fix the random seed, build
	$X=[\mathbf1\ \ x]$, and evaluate
	\verb|inv(X'*X)*X'*y|, comparing the result with the built-in
	\verb|regress|. The exercise is worth performing once for the sake of a single
	observation: $\hat\beta$ is a random vector, being a fixed linear map applied to
	the random $y$, and its realisation differs from the true $(1,3)^{\transp}$ by an
	amount that is itself the projection of the error vector $\varepsilon$,
	\[
		\hat\beta-\beta=\bigl(X^{\transp}X\bigr)^{-1}X^{\transp}\varepsilon .
	\]
	Every result in the asymptotic theory of least squares is a statement about that
	expression.
\end{eg}

\subsection{The best linear predictor}

\begin{proposition}[Population least squares]\label{prop:best-linear-predictor}
	Let $(X,Y)$ be square-integrable random variables with $\Var[X]>0$. The
	minimiser of $\E\bigl[(Y-a-bX)^2\bigr]$ over $(a,b)\in\R^2$ is
	\[
		b^\ast=\frac{\Cov(X,Y)}{\Var[X]},\qquad a^\ast=\E[Y]-b^\ast\E[X],
	\]
	and the prediction error $Y-a^\ast-b^\ast X$ has mean zero and is uncorrelated
	with $X$.
\end{proposition}

\begin{proof}
	Work in the Hilbert space $L^2$ with $\ip{U}{V}=\E[UV]$
	(\autoref{eg:ip-spaces}) and let $M\eqdef\Span\{1,X\}$, a two-dimensional,
	hence closed, subspace by \autoref{cor:subspace-complete} applied inside the
	finite-dimensional span. Minimising $\E[(Y-a-bX)^2]=\norm{Y-a-bX}^2$ is
	projecting $Y$ onto $M$, so by \autoref{prop:proj-subspace-char} the solution is
	characterised by $\E[(Y-a-bX)\cdot1]=0$ and $\E[(Y-a-bX)X]=0$. The first gives
	$a=\E[Y]-b\E[X]$; substituting into the second gives
	$\E[(Y-\E Y)X]-b\,\E[(X-\E X)X]=0$, that is $\Cov(X,Y)=b\Var[X]$.
\end{proof}

The formulas are the population counterparts of \eqref{eq:ols-solution}, and the
argument is the same projection applied in $L^2$ rather than in $\R^n$: the
sample version replaces $\E$ by the sample average, that is, replaces the $L^2$
inner product by $\ip{u}{v}=n^{-1}\sum_iu_iv_i$. Enlarging $M$ from
$\Span\{1,X\}$ to the set of all square-integrable functions of $X$ turns the
best \emph{linear} predictor into the best predictor, and
\autoref{thm:best-predictor} identifies it as $\E[Y\mid X]$.

\section{Chapter Summary}\label{sec:proj-summary}

Orthogonal complements are always closed subspaces; orthogonal sets of non-zero
vectors are independent; and orthogonality makes norms add, which is Pythagoras.
In a Hilbert space, every non-empty closed convex set contains a unique nearest
point to any given point, with existence obtained from the parallelogram law and
completeness rather than from compactness; on a subspace the characterising
variational inequality becomes the orthogonality condition
\eqref{eq:orthogonality-condition}. Hence $H=M\oplus M^{\perp}$ for every closed
subspace $M$, the projection operator is linear, idempotent, self-adjoint and
norm-reducing, and every continuous linear functional on a Hilbert space is an
inner product with a fixed vector.

In $\R^n$ the projection onto the column space of a full-column-rank $X$ is
$X(X^{\transp}X)^{-1}X^{\transp}$, and invertibility of $X^{\transp}X$ is exactly
independence of the columns. Least squares is that projection: the normal
equations are orthogonality of the residual to the regressors, the
analysis-of-variance identity is Pythagoras, $R^2$ is a squared cosine, adding
regressors cannot worsen the fit because it enlarges the subspace, and the
estimator is linear in the data. The same projection performed in $L^2$ gives the
best linear predictor, with covariance over variance replacing the sample formula.

\begin{exercise}\label{ex:proj-1}
	Let $M_1\subseteq M_2$ be closed subspaces of a Hilbert space. Show that
	$P_{M_1}P_{M_2}=P_{M_2}P_{M_1}=P_{M_1}$, and deduce that
	$\norm{y-P_{M_2}y}\leq\norm{y-P_{M_1}y}$. Give the econometric reading, and
	explain why the inequality is generally strict even when the added regressor is
	independent of $y$ in the population.
\end{exercise}

\begin{exercise}\label{ex:proj-2}
	(Frisch--Waugh--Lovell.) Partition $X=[X_1\ \ X_2]$ and let
	$Q_1=I-X_1(X_1^{\transp}X_1)^{-1}X_1^{\transp}$. Show that the sub-vector
	$\hat\beta_2$ of \eqref{eq:ols-matrix} equals the least-squares coefficient from
	regressing $Q_1y$ on $Q_1X_2$. Identify at which step
	\autoref{prop:projection-operator}(ii) and (iii) are used, and interpret the
	result as ``controlling for $X_1$''.
\end{exercise}

\begin{exercise}\label{ex:proj-3}
	Let $H=L^2([0,1])$ and $M$ the subspace of polynomials of degree at most $2$.
	Compute $P_M f$ for $f(t)=e^t$ by applying \autoref{prop:gram-schmidt} to
	$\{1,t,t^2\}$, and compare the resulting quadratic with the second-order Taylor
	polynomial of $e^t$ at $t=0$. Explain in terms of the two different
	optimality criteria why they differ.
\end{exercise}
